% !PRINTED 56
\weberpage{56}
%<<CONT

\begin{equation}
\begin{aligned}
 f'(\alpha_1)&=(\alpha_1-\alpha_2)(\alpha_1-\alpha_3)\cdots(\alpha_1-\alpha_n),\\
 f'(\alpha_2)&=(\alpha_2-\alpha_1)(\alpha_2-\alpha_3)\cdots(\alpha_2-\alpha_n),\\
 &\quad\cdots\\
 f'(\alpha_n)&=(\alpha_n-\alpha_1)(\alpha_n-\alpha_2)\cdots(\alpha_n-\alpha_{n-1}).
\end{aligned}
\tag{12}
\end{equation}

\begin{center}
{\large \S.\ 15.}

\vspace{0.5em}
{\large\gesperrt{Whole Functions of Several Variables: Forms.}}
\end{center}

Up to now we have chiefly considered whole rational functions of one variable; but we cannot always restrict ourselves to this, and indeed we have already used functions of several variables above. By a whole rational function of degree $n$ in several variables $F(x,y,z,\ldots)$ we understand a sum of terms
\[
 \sum A_{\alpha,\beta,\gamma,\ldots}x^\alpha y^\beta z^\gamma\cdots,
\]
where $\alpha,\beta,\gamma,\ldots$ are non-negative integral exponents whose sum $\alpha+\beta+\gamma+\cdots$ does not exceed the value $n$, and in at least one term actually attains it. Thus the degree is determined by the greatest value assumed by the sum $\alpha+\beta+\gamma+\cdots$.

If the sum of the exponents $\alpha+\beta+\gamma+\cdots$ has the same value in all terms, then the function is called \gesperrt{homogeneous}.

\gesperrt{A fundamental property of homogeneous functions of degree $n$ is this: if all variables are multiplied by the same factor, the result is the same as if the function were multiplied by the $n$th power}; in signs, if $t$ denotes an arbitrary variable,
\begin{equation}
 F(tx,ty,tz,\ldots)=t^nF(x,y,z,\ldots);
\tag{1}
\end{equation}
for if in the product $x^\alpha y^\beta z^\gamma\cdots$ the variables are replaced by $tx,ty,tz,\ldots$, it receives the factor
\[
 t^{\alpha+\beta+\gamma+\cdots}.
\]
If now $\alpha+\beta+\gamma+\cdots$ has in all terms the same value $n$, the factor $t^n$ can be taken out before the sum. But if the sum $\alpha+\beta+\gamma+\cdots$ has different values in the separate terms, no such common factor can be taken out, at least
%CONT>>
